\documentclass[11pt]{article}
\usepackage[a4paper,margin=27mm]{geometry}
\usepackage{amsmath,amssymb,amsthm,xcolor,booktabs,array}
\newcommand{\PROVED}{\textcolor{green!40!black}{\textbf{PROVED}}}
\newcommand{\OPEN}{\textcolor{red!70!black}{\textbf{OPEN}}}
\newcommand{\AUDIT}{\textcolor{orange!70!black}{\textbf{AUDIT}}}
\newtheorem{theorem}{Theorem}
\newtheorem{proposition}{Proposition}
\title{Source audit of the Witt--Loewy contractivity estimate (8c) (v100)}
\date{13 August 2026}

\begin{document}
\maketitle

\section{Only one orientation is needed}

Use the conventions of row (c):
\[
 (\lambda_t f)(x)=f(t^{-1}x),\qquad
 U_n=\lambda_n^{-1}=\lambda_{1/n}.
\]
On the transpose character with Mellin parameter $\rho$, $\lambda_t$ acts
by $t^\rho$.  Hence the half-Tate centered Witt operator
\begin{equation}
                    S_n:=\sqrt n\,U_n                \tag{1}
\end{equation}
acts on the corresponding simple factor by
\begin{equation}
                    \chi_\rho(S_n)=n^{1/2-\rho}.     \tag{2}
\end{equation}

\begin{theorem}[One-sided Witt contractivity closes the spectral gap ---
\PROVED]
Assume that for one integer $n\ge2$, on every simple Loewy factor of the
odd row-(c) representation,
\begin{equation}
                         \|\pi(S_n)\|\le1.           \tag{3}
\end{equation}
Then every nontrivial spectral parameter satisfies
$\Re\rho=1/2$.  Consequently it is unnecessary to prove both inequalities
in (8c); (3), together with the already proved Fourier--Poisson involution,
is enough.
\end{theorem}

\begin{proof}
Equations (2)--(3) give
$n^{1/2-\Re\rho}\le1$, hence $\Re\rho\ge1/2$.  Row (c) is invariant under
the functional-equation involution $\rho\mapsto1-\bar\rho$.  Applying the
same estimate to the mirror parameter gives
$1-\Re\rho\ge1/2$, hence $\Re\rho\le1/2$.  Both inequalities force
$\Re\rho=1/2$.
\end{proof}

This is a genuine simplification of v99: the sought norm promotion is
one-sided and can be attacked directly from the Verschiebung shift.

\section{The algebraic Witt map does not extend to the required Hilbert map}

The nuclear Dirichlet algebra in row (a) is
\[
 \mathcal C_{\mathbb R}
 =\{a:q_k(a)=\sum_{r\ge1}|a_r|r^k<\infty
       \text{ for every }k\ge0\}.
\]
The row-(b) intertwiner is algebraically
\[
 J\phi_r=\delta_r.
\]
In the half-Tate Hilbert norm of v96,
$\|\phi_r\|_{1/2}=r^{-1/2}$.

\begin{proposition}[Failure of the evident continuous promotion ---
\PROVED]
The algebraic map $J$ has no continuous extension from the half-Tate Witt
Hilbert completion to $\mathcal C_{\mathbb R}$.  The same conclusion holds
for Haran's original Hilbert norm.
\end{proposition}

\begin{proof}
For every $k\ge0$,
\[
 q_k(J\phi_r)=q_k(\delta_r)=r^k,
 \qquad
 \frac{q_k(J\phi_r)}{\|\phi_r\|_{1/2}}=r^{k+1/2}\longrightarrow\infty.
\]
Thus no defining seminorm $q_k$ is bounded by the half-Tate Hilbert norm.
For Haran's norm $\|\phi_r\|_H=\sqrt{\varphi(r)}$, take $k=1$ and use
$\varphi(r)\le r$:
\[
 \frac{q_1(\delta_r)}{\|\phi_r\|_H}
 \ge\sqrt r\longrightarrow\infty.
\]
Hence neither Hilbert completion maps continuously into the nuclear
coefficient algebra through $J$.
\end{proof}

Therefore a simple factor of row (c) is not presently a continuous Hilbert
quotient of the Witt shift.  Algebraic covariance alone does not allow the
quotient norm to be transported.

\section{Why the metric line alone cannot prove an operator bound}

\begin{proposition}[Metric/contact data do not determine operator norm ---
\PROVED]
Fix the multiplicative correspondence labels $\Gamma_n$, their metric
lines of norm $n^{-1/2}$ and their contact degrees $\Lambda(n)$.  For every
$\alpha\in\mathbb C$, the assignment
\begin{equation}
             \mathcal R_\alpha(\Gamma_n)=n^\alpha\in\operatorname{End}(\mathbb C)
                                                               \tag{4}
\end{equation}
is a strict monoidal realization and leaves the metric-line and contact
data unchanged.  Its centered operator norm is
\begin{equation}
             \|n^{-1/2}\mathcal R_\alpha(\Gamma_n)\|
             =n^{\Re\alpha-1/2}.                    \tag{5}
\end{equation}
Thus the scalar identity $\|\mathcal T_n\|^{1/2}=n^{-1/2}$ does not imply a
contractive bound on the image of $\Gamma_n$ unless the realization is
proved to be metric-enriched or dagger-contract-ive.
\end{proposition}

\begin{proof}
Multiplicativity follows from $(mn)^\alpha=m^\alpha n^\alpha$.  The
determinant line and the number $\Lambda(n)$ are separate components of the
typed row-(b) correspondence, so (4) does not alter them.  Formula (5) is
the scalar operator norm.  Taking $\Re\alpha>1/2$ violates contractivity
while preserving all the stated algebraic and scalar data.
\end{proof}

This proposition does not modify the actual row-(c) character.  It proves
the logical point needed in the audit: contractivity cannot be inferred
from the existing labelwise determinant comparison; a norm-enriched
comparison theorem is required.

\section{The obstruction already occurs inside the actual source space}

For $s\in\mathbb C$ define the Mellin functional
\begin{equation}
             \ell_s(f)=\int_0^\infty f(x)x^s\,d^*x.             \tag{5a}
\end{equation}

\begin{proposition}[Nuclear continuity does not imply critical
contractivity --- \PROVED]
For every $s\in\mathbb C$, $\ell_s$ is a continuous functional on
$\mathcal H_-=\mathcal S_{(-\infty,\infty)}$.  Its kernel is a closed
scaling-invariant hyperplane and the quotient character obeys
\begin{equation}
 \ell_s(U_nf)=n^{-s}\ell_s(f),\qquad
 \ell_s(\sqrt n\,U_nf)=n^{1/2-s}\ell_s(f).           \tag{5b}
\end{equation}
Consequently this one-dimensional nuclear quotient is contractive for
$\sqrt nU_n$ exactly when $\Re s\ge1/2$, although $\sqrt nU_n$ is unitary
on the ambient $L^2(dx)$ completion.
\end{proposition}

\begin{proof}
Choose real numbers $a<\Re s<b$.  By definition of
$\mathcal S_{(-\infty,\infty)}$, both $x^af(x)$ and $x^bf(x)$ are Schwartz
in logarithmic coordinates.  Splitting the integral at $1$ bounds (5a) by
one defining seminorm near zero and another near infinity, proving
continuity.  A change of variables gives
\[
 \ell_s(U_nf)=\int_0^\infty f(nx)x^s\,d^*x
 =n^{-s}\ell_s(f),
\]
which proves (5b).  A continuous functional has closed kernel, and (5b)
gives invariance.  The scalar norm is $n^{1/2-\Re s}$.  Finally
$\|\sqrt nU_nf\|_{L^2(dx)}=\|f\|_{L^2(dx)}$ by the same change of variables;
the apparent contradiction is absent because $\ell_s$ is not continuous
for the $L^2(dx)$ norm.
\end{proof}

By the row-(c) spectral theorem, precisely the Mellin characters attached
to the cokernel of $Z$ survive in the odd quotient.  On the initial
half-plane the identity
\[
                     \ell_s(Zf)=\zeta(s)\ell_s(f)
\]
and its source-defined continuation explain the mechanism: at a spectral
zero the functional annihilates $Z\mathcal H_\cap$ and descends to
$\mathcal H_-^0$.  Thus the existing nuclear topology deliberately retains
the very off-critical characters which ordinary $L^2$ would collapse.
Proving (3) is therefore a new boundary theorem, not a consequence of the
continuity already stated in row (c).

\section{The exact new construction which would prove (3)}

Let $M$ run over the finite Loewy subquotients of the odd object.  It is
enough to construct, directly from rows (a)--(c), Hilbert spaces
$\mathscr K_M$ and injective maps
\begin{equation}
 \Theta_M:\operatorname{gr}_LM\longrightarrow\mathscr K_M             \tag{6}
\end{equation}
such that
\begin{equation}
 \Theta_M(\sqrt n\,U_n)=V_{n,M}\Theta_M,
 \qquad \|V_{n,M}\|\le1.                            \tag{7}
\end{equation}
Then (3) follows from injectivity, the one-sided theorem gives the critical
line, and the global Hilbert assembly of v99 gives row (d).

The map $J$ is not such a $\Theta_M$ by the preceding proposition.  The
files contain no other map from the Witt Hilbert completion to the Loewy
subquotients.  Accordingly (6)--(7) are \OPEN.  Defining the norm on
$\operatorname{gr}_LM$ by pulling it back from the desired conclusion
would be circular; constructing $\Theta_M$ from a contact complex or a
Poisson boundary frame would be a valid proof.

\begin{center}
\begin{tabular}{>{\raggedright\arraybackslash}p{.57\linewidth}
                >{\centering\arraybackslash}p{.15\linewidth}
                >{\raggedright\arraybackslash}p{.18\linewidth}}
\toprule
Statement & Status & Evidence \\
\midrule
One-sided contractivity suffices & \PROVED & (1)--(3) \\
Algebraic Witt intertwiner $J$ & \PROVED & row (b) \\
$J$ extends from either Witt Hilbert completion & false & Proposition 2 \\
Metric determinant implies operator contractivity & false in this type & (4)--(5) \\
Source construction of $\Theta_M$ satisfying (6)--(7) & \OPEN & not in rows (a)--(c) \\
Row (d), condition $G$ & \OPEN & immediate after (6)--(7) \\
\bottomrule
\end{tabular}
\end{center}

\end{document}
